\section{Position and Future Work}
\label{sec:position}

\subsection{What This Enables}

Coupling detection is the foundation step of a larger program: building
an empirically-discovered, maintainable coupling model of a live system.
The push/pause/hold protocol produces a coupling scan---a directed graph
of measured step responses between subsystems, valid for the moment it
was taken. This scan is not an end in itself but the basis for:

\paragraph{Characterization.}
The probe data collected during detection already contains more
information than the binary ``coupling detected.'' Each edge carries a
measured step amplitude, a baseline noise level, and a reversibility
fraction. Extracting propagation delay, settling time, and response
curve shape from the same data requires only signal processing---no
additional probing. The scan becomes a characterized connectome:
weighted, directional, per-channel.

\paragraph{Prediction.}
If measured edges compose along graph paths, the coupling graph
becomes a forward model: ``if node $A$ changes by $\Delta$, what is
the predicted response at distant node $D$ through the chain $A \to B
\to C \to D$?'' Whether edges compose linearly or require a learned
composition function is the open question. Chain4 validation (4-node
linear topology, 3 edges) showed measured edges match planted ground
truth---but predicting a distant node's response from composed edges
is not yet tested.

\paragraph{Maintenance.}
A frozen coupling scan serves as a reference model. The system drifts;
the scan ages. Prediction error---comparing the model's expected
response against actual behavior---localizes where structure changed.
This triggers targeted rescanning of the affected neighborhood rather
than a full system rescan. The ratio of scanning to operating adapts
to structural volatility: stable systems need infrequent rescans,
volatile systems need frequent ones.

\paragraph{Cooperation.}
Once agents understand their coupling, they can act on it. Two agents
connected by a coupling edge currently find individually optimal but
jointly suboptimal configurations. The coupling graph reveals
tradeoffs invisible to independent agents. This requires no change to
the detection protocol---only a consumer of the graph it produces.

\subsection{Limitations}

\paragraph{Topology.}
The scanner discovers step-change coupling between node pairs. It does
not distinguish direct edges from transitive paths. In a branching
topology ($A \to B \to D$ and $A \to C \to D$), node $D$ responds to
$A$'s perturbation through transitive paths. The scanner reports a
step change for $(A, D)$ even though no direct edge exists.
Distinguishing direct from transitive coupling is an open problem.

\paragraph{Threshold calibration.}
The CFAR threshold uses per-sample noise (MAD of individual values)
but the detection test compares block medians. The actual false alarm
rate is lower than the nominal $P_{\text{fa}}$---the threshold is
conservative, confirmed empirically by zero false positives across all
noise levels. A median-aware threshold would improve sensitivity but
requires re-deriving the false alarm statistics.

\paragraph{Coordination scale.}
The lease-based turn-taking supports a small number of agents (tested
with 2--4). Scaling to many agents requires group-based coordination
(scoping lease tables by experiment tag) to avoid global contention.
The current global coordination table is a prototype limitation, not
a fundamental constraint.

\paragraph{Auto-calibration.}
Detection parameters---block sizes, push amplitude, cadence,
$P_{\text{fa}}$---are currently fixed. Metaheuristic search could
calibrate these per system, adapting the observation protocol to
each substrate's characteristics. This is future work.
